% TOP_PROBLEM: {"id": 30007074, "title": "Uniform log-Iitaka systems with fixed coefficients", "category_id": 5, "category": {"id": 5, "name": "algebraic_geometry", "display_name": "Algebraic Geometry", "description": "Geometric objects defined by polynomial equations.", "slug": "algebraic-geometry", "order_index": 5, "created_at": "2026-07-31T15:26:25.671Z"}, "status": "open"}
% FIRST_POSTED: 2026-10-06
% ENTRY_KIND: statement_only
% !TeX program = lualatex
\documentclass[11pt]{article}
\usepackage[margin=25mm]{geometry}
\usepackage{fontspec}
\setmainfont{Latin Modern Roman}
\usepackage{amsmath,amssymb,mathtools}
\usepackage{unicode-math}
\setmathfont{Latin Modern Math}
\usepackage{xurl,hyperref}
\hypersetup{colorlinks=true,urlcolor=blue}
\setcounter{secnumdepth}{0}
\setlength{\parindent}{0pt}
\setlength{\parskip}{5pt}
\setlength{\emergencystretch}{3em}
\newcommand{\R}{\mathbb R}
\newcommand{\N}{\mathbb N}
\newcommand{\Z}{\mathbb Z}
\newcommand{\Q}{\mathbb Q}
\newcommand{\C}{\mathbb C}
\newcommand{\D}{\mathbb D}
\newcommand{\F}{\mathbb F}
\newcommand{\sref}[2]{\hyperlink{source:#1:#2}{[S#2]}}
\newcommand{\cataloguescope}[1]{#1}
\begin{document}
\section*{Uniform log-Iitaka systems with fixed coefficients}\label{30007074}
\subsection{Definitions and mathematical statement}
\paragraph{Statement recorded in the supplied source.}
For each $d\geq 5$ and finite $\Phi \subset [0,1]\cap Q$, does one $m(d,\Phi )>0$ work for every normal integral projective lc d-fold (X,B) over any algebraically closed characteristic-zero field, with nonzero B coefficients in $\Phi $, rational Cartier $D=K_X+B$ and $\kappa (D)\geq 0$, so that $H^{0}(X,O_X(\lfloor mD\rfloor ))$ is nonzero and its ratios generate the entire Iitaka field K(D) inside k(X)?
\subsection{Short English statement}
For each $d\geq 5$ and finite $\Phi \subset [0,1]\cap Q$, does one $m(d,\Phi )>0$ work for every normal integral projective lc d-fold (X,B) over any algebraically closed characteristic-zero field, with nonzero B coefficients in $\Phi $, rational Cartier $D=K_X+B$ and $\kappa (D)\geq 0$, so that $H^{0}(X,O_X(\lfloor mD\rfloor ))$ is nonzero and its ratios generate the entire Iitaka field K(D) inside k(X)?
\subsection{Sources}
\begin{itemize}
\item[S1] Original question source or mathematical context.
\url{https://arxiv.org/abs/1005.2796}.
\item[S2] Original question source or mathematical context.
\url{https://arxiv.org/abs/2406.18233}.
\item[S3] Original question source or mathematical context.
\url{https://arxiv.org/abs/2505.05250}.
\item[S4] Original question source or mathematical context.
\url{https://arxiv.org/abs/2601.15138}.
\item[S5] Original question source or mathematical context.
\url{https://arxiv.org/abs/2604.03879}.
\item[S6] Original question source or mathematical context.
\url{https://arxiv.org/abs/2604.28085}.
\item[S7] Original question source or mathematical context.
\url{https://arxiv.org/html/2501.04232v1}.
\item[S8] Original question source or mathematical context.
\url{https://ems.press/journals/lem/articles/13750344}.
\item[S9] Original question source or mathematical context.
\url{https://www.ams.org/jams/2010-23-02/S0894-0347-09-00649-3/viewer/}.
\item[S10] Original question source or mathematical context.
\url{https://www.cambridge.org/core/journals/nagoya-mathematical-journal/article/two-nonvanishing-results-concerning-the-anticanonical-bundle/1E68F2B88B818A7CBD47B96B7D700548}.
\item[S11] OpenAI, manuscript citation page. Source of the formulation motivating this local problem.
\url{https://github.com/openai/math/blob/adc7f1241b42e322a6451854ab7e4b4c146bf78a/preprints/Uniform-log-Iitaka-fibrations-and-bounded-moduli-denominators-October-4-2026/README.md}.
\end{itemize}
\end{document}
